\section{The local-factor moment problem and the 2026 spectral letter}
\label{sec:ACspectralfollowup}

Here Q24 denotes Connes--Consani--Moscovici, \emph{On \(q\)-series and
the moment problem associated to local factors}, arXiv:2403.01247, supplied
\path{LIT_CC_q_series_moment_problem_local_factors/pqfile1.tex}.
L26 denotes Connes, \emph{The Riemann Hypothesis: Past, Present and a
Letter Through Time}, arXiv:2602.04022, supplied
\path{LIT_Connes_RH_past_present_letter/rhready.tex}.
Literal source labels are used below.  The label \texttt{spec} in Q24
is placed after the proposition environment; its statement is the
proposition with label \texttt{orthopol1}.  Neither the observed
integrality in Q24 nor the numerical agreements in L26 are treated as
unproved simply because they were first discovered experimentally:
Q24 subsequently proves its two initial ``Facts''.

\subsection{Exact inventory of the moment paper}

Q24, Proposition \texttt{momentdet}, proves determinacy of each finite
semilocal moment problem unconditionally.  Proposition
\texttt{momentssm}, equation \texttt{expth}, gives its generating
function as a sum over relatively prime semilocal pairs.  For one
prime, Proposition \texttt{expth3}, equation \texttt{normoments1},
gives geometrically convergent Lambert series in \(q=p^{-1}\).
The observed determinant integrality in Fact \texttt{keyfact} is
proved by Theorem 5.9, label \texttt{integralitymain}.  The leading
valuation \(q^{m^2}\) in Fact \texttt{keyfact2} is proved by Lemma
\texttt{keylem2}.  Theorem \texttt{eigenthm} gives the individual
eigenvalue expansions, whose coefficients belong to
\(\sqrt2\mathbb Q\); the source explicitly exhibits odd prime
denominators, so these are not asserted to be dyadic integers.
Theorem \texttt{qthm} gives characteristic-polynomial coefficients in
\(\mathbb Z[1/\sqrt2][[q]]\).  These are unconditional finite-stage
statements, not consequences of \(A\).

The following distinguishes the actual spectral limit of these Jacobi
compressions from a zeta-zero realization.

\begin{theorem}[AC-S14: the finite semilocal Jacobi limit]
\label{AC-S14}
Let \(S_f\) be a finite set of primes, \(|\alpha_p|\leq1\), and
\(\epsilon\in\{0,1\}\).  Normalize to mass one the measure
\[
 d\mu(t)=Z^{-1}
 \left|\Gamma\!\left(\frac{1/2+\epsilon+it}{2}\right)\right|^2
 \prod_{p\in S_f}|1-\alpha_pp^{-1/2+it}|^{-2}\,dt.
\]
This includes Q24 and finite collections of primitive Dirichlet local
factors.  Its moment problem is determinate, its orthogonal polynomials
are a complete basis, and their span is an operator core for
\(M_tf(t)=tf(t)\).  In that basis let \(J\) be multiplication by
\(t\), and \(J_N\) its compression to the first \(N\) coordinates,
extended by zero on the orthogonal complement.  Then
\[
 J_N\longrightarrow J\quad\hbox{in strong resolvent sense}.
\]
The \(N\) roots of the degree-\(N\) orthogonal polynomial are the
simple real eigenvalues \(x_{j,N}\) of the compression.  Successive
sets strictly interlace.  With normalized eigenvectors \(v_{j,N}\),
the positive probability measures
\[
 \mu_N=\sum_{j=1}^N|\langle e_0,v_{j,N}\rangle|^2\delta_{x_{j,N}}
\]
satisfy
\[
 \int t^k\,d\mu_N=\int t^k\,d\mu\quad(0\leq k\leq2N-1),
 \qquad \mu_N\Rightarrow\mu.
\]
Moreover the largest quadrature weight tends to zero, and every open
real interval contains a compression eigenvalue for all sufficiently
large \(N\).  All these conclusions are unconditional.  They describe
convergence to an absolutely continuous measure of full support, not
to the discrete zero measure of an \(L\)-function.
\end{theorem}
\begin{proof}
Each finite Euler factor is bounded above and below by positive
constants.  Stirling's formula gives an exponentially decaying gamma
density, with exponential moments of every rate smaller than
\(\pi/2\).  Thus all moments exist and are bounded by a constant
times \(C^{2n}(2n)!\).  Carleman's criterion proves determinacy;
this is the same argument as Q24, Proposition \texttt{momentdet}.

For completeness, polynomial density also holds in the graph norm.
Use the finite measure \(d\nu=(1+t^2)d\mu\).  If \(g\in L^2(\nu)\)
is orthogonal to every polynomial, the Fourier transform of the finite
complex measure \(g\,d\nu\) extends holomorphically to a strip about
the real axis, by Cauchy--Schwarz and the exponential moments.  Every
derivative at zero vanishes.  The identity theorem and uniqueness of
Fourier transforms of finite measures imply \(g=0\).  Hence the
polynomials are dense in \(L^2(\nu)\), which is exactly the graph
space of \(M_t\).  The same argument without \(1+t^2\) gives
ordinary completeness.

Orthogonalizing powers gives a real Jacobi matrix with strictly
positive off-diagonal entries: the measure has positive density on
all of \(\mathbb R\), so no nonzero polynomial has zero norm.  For
each finite sequence \(v\), one has \(J_Nv=Jv\) for sufficiently
large \(N\).  If \(z\notin\mathbb R\), the set of vectors
\((J-z)v\) with finite \(v\) is dense, by the core assertion.
On that set the resolvents \((J_N-z)^{-1}\) eventually agree with
\((J-z)^{-1}\).  Their norms are at most
\(|\operatorname{Im}z|^{-1}\), so they converge strongly everywhere.

The three-term recurrence identifies the monic characteristic
polynomial with the monic orthogonal polynomial; this is also Q24,
Proposition \texttt{orthopol1}.  The eigenvector recurrence implies
that its first coordinate is nonzero and determines all other
coordinates.  Thus every eigenvalue is simple and every displayed
weight is positive.  Cauchy interlacing is strict since consecutive
characteristic polynomials cannot share a root, again by the
three-term recurrence and the nonzero off-diagonal coefficients.

In a Jacobi matrix a path starting at index zero cannot exit the
first \(N\) coordinates and return in fewer than \(2N\) steps.
Expansion of the matrix powers therefore gives the asserted exact
moments through degree \(2N-1\).  The fixed second moment gives
tightness of \(\mu_N\).  Any weakly convergent subsequence has all
moments equal to those of \(\mu\): for any fixed degree use the
next higher even moment, also eventually fixed, to control the tails.
Determinacy identifies the limit, proving weak convergence of the
whole sequence.

The limiting measure has no atoms.  If some weights were at least
\(a>0\), tightness would keep their locations in a bounded set along
a subsequence.  A further subsequence converges to a point \(x\);
the portmanteau theorem applied to arbitrarily small closed
neighborhoods of \(x\) would give \(\mu(\{x\})\geq a\), a
contradiction.  Finally every open interval has positive \(\mu\)
mass, so the open-set part of that theorem gives eventually positive
\(\mu_N\) mass there.  The spectral type of \(M_t\) follows directly
from its strictly positive Lebesgue density.
\end{proof}

\begin{proposition}[AC-S15: the precise formal Witt interpretation]
\label{AC-S15}
Let \(R=\mathbb Z[1/\sqrt2]\) and use Q24's normalized determinants
\(\Delta_n(q)=D_n(q)/D_n(0)\), with \(\Delta_{-1}=1\).
The source integrality theorem gives
\(\Delta_n\in1+qR[[q]]\), and its Jacobi formula gives
\[
 \frac{a_n(q)^2}{a_n(0)^2}
   =\frac{\Delta_{n-1}(q)\Delta_{n+1}(q)}{\Delta_n(q)^2}
      \in1+qR[[q]].
\]
For the ghost convention in Q24, namely
\(g(f)=-qf'(q)/f(q)\), one has the exact identity
\[
 g\!\left(a_n^2/a_n(0)^2\right)
        =g(\Delta_{n-1})+g(\Delta_{n+1})-2g(\Delta_n).
\]
If \(\Delta_n=\det(I+\mathcal A_n(q))\), then
\[
 g(\Delta_n)=-q\operatorname{Tr}
          ((I+\mathcal A_n)^{-1}\mathcal A_n').
\]
These are identities of formal power series and give the stated
second finite difference in the additive group of the big Witt ring.
They do not assign a scheme, a motive, or an etale Frobenius to these
series.
\end{proposition}
\begin{proof}
A series with constant term one is invertible over \(R[[q]]\),
by recursive coefficient comparison.  The displayed Jacobi quotient
and the product rule for the formal logarithmic derivative give the
first two assertions.  Differentiating a determinant over formal
power series gives the last identity; its inverse exists since
\(\mathcal A_n(0)=0\).  Q24, Corollary \texttt{ansq1} and the
paragraph following it, identify this group law with Witt addition.
No geometric realization follows merely from membership in this
formal ring.
\end{proof}

\begin{proposition}[AC-S20: convergent local eigenvalue branches]
\label{AC-S20}
For either parity block of Q24's Hankel correction matrix
\(\mathcal A(q)=\mathcal D(0)^{-1}
                         (\mathcal D(q)-\mathcal D(0))\),
let \(d\) be its dimension.  There are \(d\) convergent local
eigenvalue branches
\[
 \lambda_k(q)=q^{2k+1}u_k(q),\qquad
 u_k(0)\ne0,\qquad
 u_k\in\sqrt2\mathbb Q[[q]],\quad0\leq k<d.
\]
Their convergence holds on some neighborhood of \(q=0\) depending
on the block.  This supplies a full implicit-function proof of the
expansions in Q24, Lemma \texttt{usefact2} and Theorem
\texttt{eigenthm}.  It does not give a radius uniform in the block
dimension.  These eigenvalues belong to the Hankel correction matrix,
not to the Jacobi compression of AC-S14 or a prolate operator.
\end{proposition}
\begin{proof}
The Lambert-series representation, Q24, Proposition
\texttt{similar}, is holomorphic for \(|q|<1\): for fixed matrix
indices its vector coordinates grow only polynomially in the
Lambert index, whereas the factors \(q^{2\ell+1}\) decay
geometrically on compact subdiscs.  In the monomial basis
\(B(q)=\mathcal A(q)/\sqrt2\) has rational Taylor coefficients,
since the undeformed moment matrix is \(1/\sqrt2\) times a
rational matrix and the moment corrections have rational
coefficients.

The exterior-power argument of Q24, Lemma \texttt{keylem2}, gives
\[
 e_m(q):=\operatorname{Tr}(\wedge^mB(q))
       =c_mq^{m^2}+O(q^{m^2+1}),\qquad c_m\in\mathbb Q^\times,
       \quad1\leq m\leq d,
\]
and set \(c_0=1\).  To check the nonzero coefficient, the first
\(m\) Lambert vectors, after nonzero coordinate factors, are
evaluations of Q24's polynomials \(p_0(j),\ldots,p_{m-1}(j)\)
at the \(d\) distinct indices of the chosen parity.  The polynomial
\(p_\ell\), equation \texttt{pmzg2}, has degree exactly \(\ell\).
These evaluation vectors are independent, so their Gram determinant
is strictly positive.  All scalar Lambert coefficients in the unique
lowest-degree contribution are nonzero.  This proves \(c_m\ne0\).

Write the characteristic polynomial as
\(F(q,z)=\sum_{m=0}^d(-1)^me_m(q)z^{d-m}\).
For fixed \(k\), substitute \(z=q^{2k+1}y\) and divide by
\(q^{d(2k+1)-k(k+1)}\).  Every resulting power of \(q\) is
nonnegative because
\[
 m^2+(d-m)(2k+1)-d(2k+1)+k(k+1)
       =(m-k)(m-k-1)\geq0.
\]
At \(q=0\) the resulting polynomial is
\[
 (-1)^ky^{d-k-1}(c_ky-c_{k+1}).
\]
Its nonzero root \(y=c_{k+1}/c_k\) is simple.  The analytic
implicit-function theorem gives a convergent \(y_k(q)\) there;
coefficient recursion gives rational Taylor coefficients.  The
\(d\) branches \(q^{2k+1}y_k(q)\) have distinct valuations, so
they exhaust the eigenvalues of \(B\) for small nonzero \(q\).
Multiplication by \(\sqrt2\) proves the assertion for
\(\mathcal A\).
\end{proof}

\subsection{What a direct exhaustion by critical-line local factors does}

\begin{theorem}[AC-S16: normalized semilocal measures concentrate at zero]
\label{AC-S16}
For \(X\geq2\), normalize the untwisted measure
\[
 d\mu_X(t)=Z_X^{-1}w(t)
       \prod_{p\leq X}|1-p^{-1/2+it}|^{-2}\,dt,
 \qquad w(t)=|\Gamma(1/4+it/2)|^2.
\]
Unconditionally,
\[
 \mu_X\Rightarrow\delta_0,
 \qquad \int |t|^a\,d\mu_X(t)\longrightarrow0
                  \quad\hbox{for every fixed }a>0.
\]
In particular the first Jacobi coefficient squared tends to zero.
This is an explicit result for the unscaled, mass-one normalization
of Q24's measures with all primes up to \(X\).  A discrete zeta-zero
measure is not the limit of this particular exhaustion.  It says
nothing against a differently renormalized or different geometric
construction.
\end{theorem}
\begin{proof}
The classical unconditional prime-counting estimate gives, for every
fixed \(M>0\),
\(\pi(y)=\operatorname{Li}(y)+O_M(y/\log^M y)\).
For a precise primary source, Schoenfeld, \emph{Math. Comp.} 30
(1976), Theorem 11, equations (7.1)--(7.2), p.~342, bounds
\(\psi(y)-y\) by \(y\) times a constant multiple of
\((\log y)^{1/4}\exp(-c\sqrt{\log y})\), with \(c>0\).
Prime powers of exponent at least two contribute
\(O(\sqrt y\log^2 y)\) to \(\psi-\theta\), and partial summation
of \(\theta\) gives the stated estimate for \(\pi\).
Partial summation, uniformly for real \(t\), gives
\[
 S_X(t):=\sum_{p\leq X}p^{-1/2+it}
 =\int_2^X\frac{y^{-1/2+it}}{\log y}\,dy
        +O_M\!\left((1+|t|)\frac{\sqrt X}{\log^M X}\right).
\]
Indeed the endpoint error is \(O_M(\sqrt X/\log^M X)\), and differentiation
of \(y^{-1/2+it}\) bounds the error integral by
\(C_M(1+|t|)\int_2^Xy^{-1/2}/\log^M y\,dy\), which has the
same order.
One integration by parts in the main term gives, uniformly in \(t\),
\[
 S_X(t)=\frac{X^{1/2+it}}{(1/2+it)\log X}
      +O\!\left(\frac{\sqrt X}{\log^2X}
             +(1+|t|)\frac{\sqrt X}{\log^4X}\right),
\]
where \(M=4\) has been chosen.
Thus, for each fixed \(\eta>0\), uniformly on
\(\eta\leq|t|\leq\log X\),
\[
 \sum_{p\leq X}p^{-1/2}(1-\cos(t\log p))
     =S_X(0)-\operatorname{Re}S_X(t)
     \geq c_\eta\frac{\sqrt X}{\log X}
\]
for sufficiently large \(X\).  The positive main coefficient is
\(2-1/\sqrt{1/4+\eta^2}\).

Divide the product by its value at zero and set
\[
 \rho_X(t)=\prod_{p\leq X}
       \frac{(1-p^{-1/2})^2}{|1-p^{-1/2+it}|^2},
 \qquad D_X=\int_{\mathbb R}w(t)\rho_X(t)\,dt.
\]
Then \(0<\rho_X\leq1\) and \(d\mu_X=D_X^{-1}w\rho_X\,dt\).
Writing \(r=p^{-1/2}\) and
\(u=2r(1-\cos\theta)/(1-r)^2\), the inequalities
\(u/(1+u)\leq\log(1+u)\leq u\) imply
\[
 \log\rho_X(t)
 \leq-\tfrac12\sum_{p\leq X}p^{-1/2}(1-\cos(t\log p)),
\]
because \((1+r)^2\leq4\).  Consequently
\[
 \rho_X(t)\leq\exp(-c_\eta\sqrt X/\log X)
        \quad(\eta\leq|t|\leq\log X),
\]
after decreasing \(c_\eta>0\) if necessary.

On the other hand, \(r\leq2^{-1/2}\) and
\(1-\cos\theta\leq\theta^2/2\) give
\[
 \log\rho_X(t)\geq
     -C t^2\sum_{p\leq X}p^{-1/2}(\log p)^2
     \geq-C' t^2\sqrt X\log X.
\]
The last bound uses \(S_X(0)=O(\sqrt X/\log X)\).
For \(\delta_X=X^{-1/4}(\log X)^{-1/2}\), this proves
\(\rho_X(t)\geq c>0\) when \(|t|\leq\delta_X\).
Since \(w\) is continuous and positive at zero,
\[
 D_X\geq c'\delta_X.
\]

It remains to control all heights, not just bounded \(t\).
For fixed \(a\geq0\), Stirling's formula yields
\[
 \int_{|t|\geq T}|t|^a w(t)\,dt
       \leq C_a(1+T)^a e^{-\pi T/2}\quad(T\geq1).
\]
Using \(\rho_X\leq1\), the contribution from \(|t|\geq\log X\),
after division by \(D_X\), is at most
\[
 C_a X^{1/4-\pi/2}(\log X)^{a+1/2}=o(1).
\]
The contribution from \(\eta\leq|t|\leq\log X\) is at most
\(C_a\delta_X^{-1}\exp(-c_\eta\sqrt X/\log X)=o(1)\).
Thus the mass outside every fixed neighborhood of zero tends to zero.
For \(a>0\), the remaining contribution from \(|t|<\eta\) is
at most \(\eta^a\), so all the stated absolute moments tend to zero.
The measure is even, and the first Jacobi coefficient squared is its
second moment.  This proves the final assertion.
\end{proof}

Theorem A supplies an alternative sufficient estimate:
Schoenfeld's Corollary 1, equation (6.18), pp.~339--340, gives
\(\pi(y)-\operatorname{Li}(y)=O(\sqrt y\log y)\), and then the
error in the first partial-summation formula is
\(O((1+|t|)\log^2X)\).  The proof above shows that this stronger
input is unnecessary for AC-S16.

The theorem concerns the untwisted prime exhaustion only.  The
unit-modulus phases of a nontrivial character need a separate
analysis when the prime set grows.  Finite sets, including all such
phases, are covered by AC-S14.  Neither Q24's formal integrality nor
RH by itself supplies a uniform-in-degree estimate for those
characteristic polynomials or for a semilocal prolate ground state.

\begin{theorem}[AC-S21: the central scaling has a Gaussian limit]
\label{AC-S21}
For the untwisted measures of AC-S16 put
\[
 C_X=\sum_{p\leq X}
       \frac{p^{-1/2}(\log p)^2}{(1-p^{-1/2})^2},
 \qquad s_X=\sqrt{2C_X}.
\]
Then \(C_X\sim2\sqrt X\log X\).  Unconditionally, the image
\(\nu_X\) of \(\mu_X\) under \(t\mapsto s_Xt\) converges in
total variation to the standard Gaussian probability measure, and
\[
 \int |u|^a\,d\nu_X(u)\longrightarrow
 \frac1{\sqrt{2\pi}}\int_{\mathbb R}|u|^ae^{-u^2/2}\,du
 \qquad(a>0).
\]
Let \(a_n(X)>0\) be the off-diagonal Jacobi coefficient joining
degrees \(n\) and \(n+1\) for \(\mu_X\).  For every fixed \(n\geq0\),
\[
 s_Xa_n(X)\longrightarrow\sqrt{n+1},
 \qquad
 a_n(X)^2\sim\frac{n+1}{4\sqrt X\log X}.
\]
On the common sequence space supplied by orthonormal polynomials,
the self-adjoint operators \(s_XJ_X\) converge in strong resolvent
sense to the Gaussian Jacobi operator, with coefficients
\(\sqrt{n+1}\) and zero diagonal.  This identifies the limit of this
specific central rescaling; it does not identify a zeta-zero
operator or a semilocal prolate operator.
\end{theorem}
\begin{proof}
Write \(r_p=p^{-1/2}\), and retain \(\rho_X\) from AC-S16.
The difference between \(C_X\) and
\(\sum_{p\leq X}r_p(\log p)^2\) is
\[
 O\!\left(\sum_{p\leq X}p^{-1}(\log p)^2\right)
       =O(\log^3X).
\]
The last estimate follows even by summing over all positive integers.
Partial summation with the unconditional prime-counting estimate
used in AC-S16 gives
\[
 \sum_{p\leq X}r_p(\log p)^2
  =\int_2^X y^{-1/2}\log y\,dy
           +O(\sqrt X/\log^2X)
  =2\sqrt X\log X-4\sqrt X+O(\sqrt X/\log^2X).
\]
This proves the asserted asymptotic for \(C_X\).

For \(u\) in a fixed compact set, set \(t=u/s_X\).
Since \(|t|\log X\to0\), Taylor expansion of each logarithm gives,
uniformly on that set,
\[
 \log\rho_X(t)
  =-C_Xt^2+
     O\!\left(t^4\sum_{p\leq X}
                \frac{r_p(\log p)^4}{(1-r_p)^2}\right)
  =-\frac{u^2}{2}
       +O\!\left(\frac{u^4\log^2X}{C_X}\right).
\]
The error is valid with an absolute constant: the coefficients
\(r_p/(1-r_p)^2\) are bounded, so the quadratic remainder of
\(\log(1+v)\) is bounded by a constant times the same coefficient
times the fourth power of \(t\log p\).
Consequently \(\rho_X(u/s_X)\to e^{-u^2/2}\).
Moreover, for \(|t|\leq1/\log X\), the inequalities from AC-S16
and \(1-\cos(t\log p)\geq c(t\log p)^2\) imply
\[
 \rho_X(u/s_X)\leq e^{-c'u^2}
       \quad\bigl(|u|\leq s_X/\log X\bigr),
 \tag{SF.1}
\]
because \(C_X\) and \(\sum r_p(\log p)^2\) are comparable.

We must control the complement of this growing Gaussian region.
For every fixed \(M>0\), partial summation using the prime-counting
remainder with logarithmic power six gives, uniformly for
\(|t|\leq M\log X\),
\[
 \sum_{p\leq X}r_p(1-\cos(t\log p))
 =\int_2^X
       \frac{1-\cos(t\log y)}{\sqrt y\log y}\,dy
       +O_M(\sqrt X/\log^5X).
\]
Set \(L=\log X\).  On \(Xe^{-1}\leq y\leq X\), write
\(y=Xe^{-v}\), \(0\leq v\leq1\).  The positive integral is at least
\[
 c\,\frac{\sqrt X}{L}
       \int_0^1\bigl(1-\cos(tL-tv)\bigr)\,dv.
\]
For all real \(a,t\),
\[
 \int_0^1(1-\cos(a-tv))\,dv
 \geq1-\left|\frac{\sin(t/2)}{t/2}\right|
 \geq c_0\min(t^2,1).
\]
At \(t=0\) use continuity; the final bound follows from the
Taylor expansion near zero, strict inequality away from zero on
compact sets, and the bound \(2/|t|\) for large \(|t|\).
It follows that, for sufficiently large \(X\),
\[
 \rho_X(t)\leq
       \exp(-c_M\sqrt X/\log^3X)
 \quad\left(\frac1{\log X}\leq |t|\leq M\log X\right).
 \tag{SF.2}
\]
Here the error \(O_M(\sqrt X/\log^5X)\) is smaller than the
uniform lower bound \(c\sqrt X/\log^3X\).

For a fixed \(a\geq0\), consider
\[
 I_{X,a}=\int_{\mathbb R}|u|^a
             w(u/s_X)\rho_X(u/s_X)\,du.
\]
On \(|u|\leq s_X/\log X\), (SF.1), boundedness of \(w\), and
dominated convergence give the limit
\(w(0)\int|u|^ae^{-u^2/2}\,du\).
The intermediate region in (SF.2) contributes at most
\[
 s_X^{a+1}e^{-c_M\sqrt X/\log^3X}
                   \int_{\mathbb R}|t|^aw(t)\,dt=o(1).
\]
On \(|t|\geq M\log X\), use \(\rho_X\leq1\) and the gamma tail
estimate in AC-S16.  The resulting bound is
\[
 O_a\!\left(
 X^{(a+1)/4-\pi M/2}
 (\log X)^{(a+1)/2+a}\right)=o(1)
\]
upon choosing \(M>(a+1)/(2\pi)\).
Thus \(I_{X,a}\) has the stated Gaussian limit for every \(a\geq0\).
The case \(a=0\) gives the normalizing constant.  Dividing by it
proves all absolute-moment limits.  The normalized densities also
converge pointwise to the Gaussian density and have integral one.
Since the integrals of their pointwise minima with that density
converge to one by dominated convergence, their \(L^1\) distance
tends to zero.  This proves total variation convergence.

Every fixed Hankel moment matrix for \(\nu_X\) now converges to
the corresponding positive definite Gaussian moment matrix.
Gram--Schmidt orthogonalization and its recurrence coefficients
depend continuously on these finitely many moments.  The standard
Gaussian orthonormal Hermite polynomials have off-diagonal
coefficients \(\sqrt{n+1}\), proving the two Jacobi assertions.
Both measures are even, so all diagonal coefficients vanish.

By AC-S14 each finite-prime Jacobi operator is self-adjoint on the
closure of the finite sequences, and those sequences also form a
core for the Gaussian operator \(J_{\rm G}\), by the same
exponential-moment proof.  For each finite sequence \(v\),
\(s_XJ_Xv\to J_{\rm G}v\).  For \(z\notin\mathbb R\),
\[
 (s_XJ_X-z)^{-1}(J_{\rm G}-z)v-v
  =(s_XJ_X-z)^{-1}(J_{\rm G}-s_XJ_X)v\longrightarrow0.
\]
The uniform resolvent bound and density of
\((J_{\rm G}-z)\) applied to the finite-sequence core prove strong
resolvent convergence on the full sequence space.
\end{proof}

\subsection{The exact contribution of A to the letter}

L26's Theorem \texttt{mainintro}, in section \texttt{sectvW}, assumes
a simple isolated minimum with even eigenfunction.  It does not prove
those hypotheses for the Weil form.  Its section \texttt{miss}
explicitly retains both that issue and the comparison of the
prolate trial function with the actual ground eigenfunction.  Its
numerical sections, the table in the letter, and Figure
\texttt{fpro1} report approximations rather than error-certified
limits.  AC-S03--AC-S07 and AC-S13 apply without changing that status.
The scalar correction in AC-S04 also applies to L26, equation
\texttt{ku}, which uses the same \(\pi/2\) coefficient.

The letter's conditional heat theorem, label \texttt{heatkernelintro},
equation \texttt{asymp}, applies to the explicit Hilbert operator of
AC-S08 once A is supplied.  Its full expansion is a cited theorem;
the following first two terms can also be obtained directly and
serve to check the counting and multiplicity normalization.

\begin{proposition}[AC-S17: heat trace of the zero-built realization]
\label{AC-S17}
For the operator \(D\) of AC-S08, counting both signs and all zero
multiplicities,
\[
 \operatorname{Tr}(e^{-tD^2})
 =\frac{\log(1/t)}{4\sqrt\pi\sqrt t}
 -\frac{\log(4\pi)+\gamma_{\rm E}/2}{2\sqrt\pi\sqrt t}
 +O(\log(1/t))\qquad(t\downarrow0).
\]
In particular its heat operators are trace class.  This is a theorem
about the zero-built operator; the same first terms for a different
operator do not establish equality of its spectrum with the zeros.
\end{proposition}
\begin{proof}
Let \(N(u)\) count positive ordinates with multiplicity.  The
Riemann--von Mangoldt formula gives
\(N(u)=F(u)+O(\log(2+u))\), where
\(F(u)=u(\log(u/(2\pi))-1)/(2\pi)\).  Since
\(N(u)=O(u\log(2+u))\), Stieltjes integration gives a convergent
trace \(2\int_0^\infty e^{-tu^2}\,dN(u)\).
The remainder contributes \(O(\log(1/t))\), by integration by parts
and the substitution \(v=\sqrt t\,u\).
The main term is
\[
 \frac1\pi\int_0^\infty e^{-tu^2}\log(u/(2\pi))\,du.
\]
Differentiating
\(\int_0^\infty e^{-tu^2}u^{s-1}du
   =\tfrac12t^{-s/2}\Gamma(s/2)\)
at \(s=1\), and using
\(\Gamma'(1/2)/\Gamma(1/2)=-\gamma_{\rm E}-2\log2\),
gives exactly the two displayed coefficients.
\end{proof}

\subsection{A proved large-window subspace limit without simple-even hypotheses}

\begin{theorem}[AC-S18: saturation by low-energy spectral subspaces]
\label{AC-S18}
Let \(A_L\) be the closed Weil operator and put \(\lambda=e^L\).
Let \(P_\lambda\) be its spectral projection for
\([0,e^{-\pi\lambda^2}]\), acting on \(L^2([-L,L])\), extended
as zero on its orthogonal complement in \(L^2(\mathbb R)\).
Under A these are finite-rank orthogonal projections and
\[
 P_\lambda\longrightarrow I\quad\hbox{strongly on }L^2(\mathbb R),
 \qquad \operatorname{rank}P_\lambda\longrightarrow\infty.
\]
For the function \(K\) of AC-S04 and every fixed \(j\geq0\),
\(b\geq0\),
\[
 \|P_\lambda K^{(j)}-K^{(j)}\|_{L^1(e^{b|x|}dx)}
  =O_{j,b}\!\left(
    \lambda^{b+9/2+2j}\sqrt{\log\lambda}\,
                         e^{-\pi\lambda^2/2}\right).
\]
Their Fourier transforms therefore converge uniformly on
\(|\operatorname{Im}z|\leq b\) to \((iz)^j\Xi(z)\).
The vectors \(P_\lambda K\) are even and belong to a proved
low-energy spectral subspace.  They need not be individual
eigenvectors, and the theorem does not imply a real-zero property
for their transforms before taking the limit.
\end{theorem}
\begin{proof}
Compact resolvent gives finite rank.  Put
\(f_{j,L}=1_{[-L,L]}K^{(j)}\).  The proof of AC-S05 gives
\[
 q_L(f_{j,L})\leq C_j\lambda^{9+4j}e^{-2\pi\lambda^2}.
\]
The spectral theorem and nonnegativity give
\[
 \|(I-P_\lambda)f_{j,L}\|_2^2
    \leq e^{\pi\lambda^2}q_L(f_{j,L})
    \leq C_j\lambda^{9+4j}e^{-\pi\lambda^2}.
\]
Here \(P_\lambda K^{(j)}=P_\lambda f_{j,L}\).
Weighted Cauchy--Schwarz on the window, using
\(\int_{-L}^Le^{2b|x|}dx\leq2L\lambda^{2b}\), and the
double-exponential tails of AC-S04 prove the stated weighted bound.
They also prove ordinary \(L^2\) convergence on each \(K^{(j)}\).

The span of these derivatives is dense in \(L^2(\mathbb R)\).
Indeed, Fourier transformation reduces the assertion to density of
the polynomial multiples of \(\Xi(t)\).  The alternating-series
formula gives \(\zeta(1/2+it)=O(1+|t|)\): its partial-sum function
is bounded, and \(|1-2^{1/2-it}|\geq\sqrt2-1\).
Together with Stirling's formula this shows that \(|\Xi(t)|^2dt\)
has exponential moments of every rate smaller than \(\pi/2\).
If \(g\in L^2(dt)\) is orthogonal to all polynomial multiples of
\(\Xi\), the Fourier transform of
\(g(t)\overline{\Xi(t)}dt\) is holomorphic in a strip about the
real axis and all its derivatives at zero vanish.  Fourier uniqueness
implies \(g\overline\Xi=0\) almost everywhere.  Since a nonzero
entire function has only a discrete real zero set, \(g=0\).
This proves density.  The projections have norm at most one, so
convergence on this dense span gives strong convergence everywhere.

For every fixed \(r\), AC-S05 gives
\(a_r(L)\leq C_r\lambda^{4r+5}e^{-2\pi\lambda^2}
 <e^{-\pi\lambda^2}\) for all sufficiently large \(\lambda\).
Hence the rank is eventually at least \(r\).  Reflection commutes
with the Weil form and its spectral projections, so \(P_\lambda K\)
is even.  The Fourier limit follows from the weighted norm bound.
\end{proof}

\begin{corollary}[AC-S19: finite-window closedness is essential]
\label{AC-S19}
Under A the nonnegative Weil form on \(C_c^\infty(\mathbb R)\),
viewed in unweighted \(L^2(\mathbb R)\), is not closable.  The same
is true on the union of the finite-window form domains.  This does
not contradict the closedness of each fixed-window form.
\end{corollary}
\begin{proof}
The forms agree on nested domains: their prime-side distributions
agree on compact tests, and passage through the common smooth form
cores gives the same sampling identity
\(q(f)=\sum_\gamma m_\gamma|\widehat f(\gamma)|^2\)
on the union.  This defines a nonnegative form there.
Fix a nonzero \(g\in C_c^\infty(\mathbb R)\).  AC-S03 gives
\(q(g)>0\).  By AC-S18,
\(v_\lambda=P_\lambda g\to g\) in \(L^2\), whereas
\[
 q(v_\lambda)\leq e^{-\pi\lambda^2}\|g\|_2^2\longrightarrow0.
\]
Thus \(u_\lambda=g-v_\lambda\to0\) in \(L^2\), and
\[
 q(u_\lambda-u_\mu)
       \leq2q(v_\lambda)+2q(v_\mu)\longrightarrow0.
\]
The reverse triangle inequality for the seminorm \(q^{1/2}\)
gives \(q(u_\lambda)\to q(g)>0\), contrary to the criterion for
closability.  This proves the union-domain assertion.

To obtain the assertion on \(C_c^\infty\) itself, choose a sequence
\(\lambda_n\to\infty\).  In each finite window approximate
\(v_{\lambda_n}\) by a compactly supported smooth \(h_n\) in both
\(L^2\) and the form norm, with errors tending to zero.  This is
legitimate by the smooth form-core property, including the endpoint
smoothing estimate proved in AC-S05.  Then \(h_n\to g\) and
\(q(h_n)\to0\).  The smooth compact tests \(g-h_n\) violate
closability by the same inequalities.
\end{proof}

AC-S18 is a genuine large-window convergence theorem available from
A without proving simplicity or evenness of the minimum.  It also
shows exactly why that convergence is weaker than L26's desired
ground-state convergence: the low-energy subspaces approximate every
\(L^2\) vector, so their near-null property does not single out the
theta kernel.  Selecting the bottom eigenline still requires a
quantitative separation and an eigenvector comparison, such as the
gate in AC-S07.  No logical independence from A is asserted.
